\section{Why the Un-generated Cannot Be Handled}

\label{sec:ungenerated}

The previous section established that physical theory consistently describes
the domain after generation has occurred.

This naturally raises a question.

\medskip
\begin{quote}
Why can the domain prior to generation not be described?
\end{quote}
\medskip

One might initially assume that this domain is simply unknown,
and that a more advanced theory will eventually describe it.

However, the issue here is not merely one of incomplete knowledge.

It concerns the form of description itself.

Physical theory presupposes several basic structures:
time, space, states, and observable quantities.
These are taken as given,
and laws are defined on the basis of them.

The \ungenerated{} domain fits none of these presuppositions.

If time has not yet come to hold,
the question of ``when'' has no meaning.
If space has not yet come to hold,
the question of ``where'' does not arise.
If states are not defined,
quantities cannot be assigned.
If observable quantities do not exist,
theory has no object.

From this, a crucial point follows.

\medskip
\begin{quote}
The \ungenerated{} is not merely unknown.\\
It is not something that can, in principle, be described.
\end{quote}
\medskip

It does not fit the form of description itself.

No matter how precise a theory becomes,
as long as it presupposes structures such as time and space,
it cannot reach this domain.

This is not a limitation of current theory,
but a constraint imposed by the conditions under which theory holds.

\medskip
\begin{quote}
Because theory presupposes these structures,\\
it cannot extend beyond them.
\end{quote}
\medskip

At this point,
the significance of divergence becomes clearer.

Theory cannot describe the \ungenerated{} domain directly.
But it can approach its boundary.

As it does so,
quantities cease to be defined,
and divergence appears.

This is not accidental.

It indicates that a discontinuity exists
between the domain that can be described
and the domain that cannot.

However,
this boundary is not a sharp line.

\medskip
\begin{quote}
It is a transition band of finite thickness.
\end{quote}
\medskip

Physical theories always operate within finite precision and scale.
The \generated{} and the \ungenerated{}
cannot be cleanly separated.

Between them exists a region
in which the descriptive framework gradually loses validity.

Within this region,
theory does not immediately fail,
but becomes unstable,
and eventually produces divergence.

From this perspective,
the following statement becomes natural.

\medskip
\begin{quote}
Divergence is not reaching the \ungenerated{} domain.\\
It is contact with what lies just before it.
\end{quote}
\medskip

This contact does not allow direct description.
But it cannot be denied.

Theory approaches, responds, and then breaks down.
The manner of this breakdown itself
is the only available indication of the boundary.

Thus,
the \ungenerated{} domain does not appear within theory directly.
It appears only through its boundary.

The next section defines this boundary explicitly.

It is not introduced as an abstract limit,
but as a structure with physical meaning:
the \DH{}.
